% =============================================================================
% SECTION: Expanded Historical Context (English)
% From the rise of the houses to 1507 and beyond
% Designed for public history use alongside OpenHistoryMap
%
% @ohm:layer "Genoa 1507 — La Superba and the Italian Wars"
% @ohm:timespan "1100/1512"
% @ohm:language "en"
% =============================================================================

\chapter{The Long History of Genoa}
\label{ch:history-expanded}

\begin{quote}
\itshape
This chapter is more than background for the game. It is the raw material
for a \textbf{public history event}: twenty players embodying the great
Genoese families of 1507 are also exploring, through fiction, one of the
most acute crises of Renaissance Italian history. The choices they make
tonight mirror real ones --- and the real outcome is documented.
\end{quote}

% =============================================================================
\section{Origins: from the Commune to the Republic (1099--1339)}
\label{sec:origins}
% =============================================================================

% @ohm:Event { "id": "genova:fondazione-comune", "date": "1099",
%   "label": "Foundation of the Commune of Genoa",
%   "type": "ohm:PoliticalEvent", "place": "genova:CittaDiGenova" }

Genoa becomes a \textit{comune} in \textbf{1099}, during the First Crusade:
Genoese merchants finance the crusading fleets in exchange for commercial
quarters in Levantine cities. This is the Genoese model: not territorial
conquest, but \textbf{control of trade nodes}. Alexandria, Constantinople,
Caffa in Crimea, Ceuta --- wherever there is a strategic port, there is a
Genoese \textit{fondaco} (trading house).

\begin{multicols}{2}

\subsection{Guelphs and Ghibellines}

% @ohm:Actor { "id": "genova:fazione-guelfa", "label": "Guelph Faction",
%   "type": "ohm:PoliticalFaction", "timespan": "1200/1528" }
% @ohm:Actor { "id": "genova:fazione-ghibellina", "label": "Ghibelline Faction",
%   "type": "ohm:PoliticalFaction", "timespan": "1200/1528" }

The division between \idxbold{Guelphs} and \idxbold{Ghibellines} in Genoa
is not merely ideological --- it is structural. The \textbf{Fieschi} lead
the Guelphs (loyal to the Pope); the \textbf{Doria} and \textbf{Spinola}
lead the Ghibellines (loyal to the Emperor). The \textbf{Grimaldi},
nominally Guelph but constitutionally pragmatic, shift allegiance as needed.

This division generates a coup cycle roughly every decade: when one faction
gains the upper hand, the other is exiled. The exiles return. The cycle
restarts. Genoa accumulates wealth despite --- or perhaps \textit{because of}
--- this permanent instability.

\columnbreak

\subsection{The Bank of Saint George}

% @ohm:Actor { "id": "genova:banco-san-giorgio",
%   "label": "Bank of Saint George (Banco di San Giorgio)",
%   "type": "ohm:Institution",
%   "timespan": "1407/1805",
%   "place": "genova:PalazzoDiSanGiorgio" }
% @ohm:Place { "id": "genova:PalazzoDiSanGiorgio",
%   "label": "Palazzo di San Giorgio",
%   "geo:lat": 44.4088, "geo:long": 8.9269,
%   "ohm:accuracy": "high" }

Founded in \textbf{1407}, the \idxbold{Bank of Saint George} (\textit{Banco
di San Giorgio}) is arguably the first modern public bank in Europe. It
manages Genoa's public debt, controls overseas colonies, and --- crucially
--- functions as \textbf{a state within a state}: its own laws, its own
prisons, its own diplomatic agents. During the 1507 occupation it remains
formally neutral. Whoever controls Genoa must deal with the Bank on the
Bank's terms, not the other way around.

\end{multicols}

% =============================================================================
\section{The Visconti Period and the Age of Crisis (1339--1453)}
\label{sec:crisis}
% =============================================================================

% @ohm:Event { "id": "genova:primo-doge", "date": "1339",
%   "label": "Simone Boccanegra first Doge of Genoa",
%   "type": "ohm:PoliticalEvent" }

In \textbf{1339} Genoa establishes the \textbf{Dogato} --- an attempt to
escape the factional cycle by electing a life-term doge. The first is
Simone Boccanegra (the story that Verdi will dramatise four centuries later).
It does not work: doges are assassinated, exiled, or deposed with clockwork
regularity.

\begin{multicols}{2}

Between 1353 and 1421, Genoa voluntarily places itself under the protection
of the \textbf{Visconti of Milan} three separate times --- each time to
escape an internal crisis that local factions cannot resolve. This pattern
--- \textit{ceding sovereignty to an outside power in exchange for order}
--- is precisely what produced the French occupation of 1499.

\columnbreak

\begin{rulebox}[SharedBrown]{Note for Players}
Your 1507 character carries this institutional memory. Genoa has
\textit{already chosen} to cede sovereignty before. The question
is not \textit{whether} it is possible --- it is \textit{on what terms}
and \textit{with what guarantees}. This is why some families negotiate
with the French even while supporting the revolt.
\end{rulebox}

\end{multicols}

% =============================================================================
\section{Apex and Decline: 1453--1499}
\label{sec:decline}
% =============================================================================

% @ohm:Event { "id": "genova:caduta-caffa", "date": "1475",
%   "label": "Fall of Caffa to the Ottomans",
%   "type": "ohm:MilitaryEvent", "place": "ohm:Caffa" }

The fall of \textbf{Constantinople} (1453) and then \textbf{Caffa} (1475)
to the Ottomans decapitates the Genoese commercial system in the eastern
Mediterranean. In a decade, Genoa loses what three centuries built.
The economic void is severe --- and opens the door to foreign creditors
and political manipulation.

\begin{multicols}{2}

\subsection{The Italian Wars Begin}

% @ohm:Event { "id": "genova:invasione-carlo-viii", "date": "1494",
%   "label": "Charles VIII invades Italy",
%   "type": "ohm:MilitaryEvent" }

In \textbf{1494} Charles~VIII of France crosses the Alps with 25,000 men
--- the largest army Italy has seen in centuries. He conquers Naples in
six months. The peninsular balance, built through the Peace of Lodi (1454)
and maintained by Lorenzo de' Medici, collapses in months. The
\idxbold{Italian Wars} (1494--1559) will turn the peninsula into a
battlefield between France, Spain, and the Empire for sixty years.

\columnbreak

\subsection{Genoa as Pawn}

Genoa is geographically critical: whoever controls northern Italy must
control its port. The Sforza of Milan --- historical allies of several
Genoese families --- are expelled by Louis~XII in \textbf{1499}.
With them falls the last buffer against direct French occupation.

% @ohm:Event { "id": "genova:occupazione-francese", "date": "1499",
%   "label": "Louis XII occupies the Duchy of Milan and Genoa",
%   "type": "ohm:MilitaryEvent",
%   "place": "genova:CittaDiGenova" }

\end{multicols}

% =============================================================================
\section{The French Occupation (1499--1507)}
\label{sec:occupation}
% =============================================================================

% @ohm:Actor { "id": "genova:luigi-xii",
%   "label": "Louis XII, King of France",
%   "type": "ohm:Person", "timespan": "1462/1515" }
% @ohm:Actor { "id": "genova:trivulzio",
%   "label": "Gian Giacomo Trivulzio",
%   "type": "ohm:Person", "timespan": "1441/1518",
%   "ohm:role": "Governor of Genoa 1499--1500, 1507" }

\begin{multicols}{2}

Under Louis~XII, Genoa is not annexed to France --- it is treated as
a \textbf{personal fief of the king}. The distinction matters: Genoa's
formal institutions (the Doge, the Bank of Saint George, the family
councils) survive, but without real power. The French governor ---
first \textbf{Trivulzio}, then others --- maintains the garrison in the
castle of \textbf{Castelletto}, which dominates the city from above.

\subsubsection{The Tax Burden}

French taxation is systematic and heavy. Louis~XII needs to finance his
Italian campaigns, and Genoa is one of his primary sources. Customs duties
are re-tariffed, commercial monopolies redistributed in favour of French
merchants, and the historical exemptions of Genoese families cancelled by
decree.

\columnbreak

\subsubsection{The Garrison}

% @ohm:Place { "id": "genova:Castelletto",
%   "label": "Castle of Castelletto",
%   "type": "ohm:MilitaryStructure",
%   "geo:lat": 44.4120, "geo:long": 8.9275,
%   "ohm:accuracy": "approximate",
%   "ohm:note": "Demolished 1536. Position is approximate." }

The \idxbold{Castelletto} is the physical symbol of occupation: a medieval
fortress on the hill overlooking the harbour, reinforced by the French and
permanently garrisoned. It is virtually impregnable from inside the city.
This is why the 1507 revolt does not plan to \textit{capture} it --- but
to \textit{isolate} it.

\begin{rulebox}[SharedBrown]{Geographic Note --- OHM}
The Castelletto was definitively demolished in 1536. Its position is
reconstructible from sixteenth-century cartographic sources.
Suggested OHM layer: \texttt{genova1507:military\_structures}.
\end{rulebox}

\end{multicols}

% =============================================================================
\section{The 1507 Revolt: Precise Chronology}
\label{sec:revolt-chronology}
% =============================================================================

% @ohm:Event { "id": "genova:rivolta-1507",
%   "label": "Revolt of Genoa against France",
%   "type": "ohm:MilitaryEvent",
%   "timespan": "1507-01/1507-04",
%   "place": "genova:CittaDiGenova" }

This is the real chronology of events that your game simulates and
compresses into a single night. It is worth knowing --- both for players
who want to anchor the fiction in history, and for understanding
\textit{why} the characters' choices carried real weight.

\begin{center}
\rowcolors{2}{LightBG}{white}
\begin{tabularx}{\textwidth}{L{3.2cm} X}
  \toprule
  \textbf{Date} & \textbf{Event} \\
  \midrule
  \textbf{Summer 1506} &
    % @ohm:Event { "id":"genova:tensioni-1506","date":"1506-07","label":"Fiscal tensions peak in Genoa","type":"ohm:SocialEvent" }
    Fiscal tensions reach a breaking point. Harbour merchants refuse to
    pay new commodity taxes. Small clashes between commoners and French guards. \\

  \textbf{January 1507} &
    % @ohm:Event { "id":"genova:primo-insurrezione","date":"1507-01","label":"First popular uprising in the caruggi","type":"ohm:MilitaryEvent" }
    First popular uprising in the \textit{caruggi}. The French respond
    with arrests and reprisals. The noble families watch --- they do not
    yet intervene. \\

  \textbf{February 1507} &
    The Genoese families begin to coordinate --- not publicly, but through
    intermediaries. Governor Trivulzio requests reinforcements from
    Louis~XII. \\

  \textbf{April 1507} &
    % @ohm:Event { "id":"genova:rivolta-principale","date":"1507-04",
    %   "label":"Main revolt --- noble families openly take sides",
    %   "type":"ohm:MilitaryEvent" }
    The revolt becomes open. The noble families declare themselves.
    The bells of San Lorenzo ring. Trivulzio retreats to the Castelletto.
    Genoa is de facto free for several weeks. \\

  \textbf{May 1507} &
    % @ohm:Event { "id":"genova:reconquista-francese","date":"1507-05",
    %   "label":"Louis XII retakes Genoa","type":"ohm:MilitaryEvent" }
    Louis~XII personally descends into Italy with an army. Genoa is
    retaken. The reprisals are severe but selective --- the king needs
    the Genoese merchants to finance his campaigns. \\

  \textbf{June 1507} &
    % @ohm:Event { "id":"genova:amnistia","date":"1507-06",
    %   "label":"Partial amnesty and new governance agreements","type":"ohm:DiplomaticEvent" }
    Partial amnesty. New governance agreements: Genoa gains some fiscal
    concessions in exchange for formal submission. Families that had
    negotiated --- such as the Grimaldi --- obtain better terms. \\
  \bottomrule
\end{tabularx}
\end{center}

% =============================================================================
\section{After 1507: What Really Happened}
\label{sec:aftermath}
% =============================================================================

\begin{multicols}{2}

This section breaks the fourth wall of the fiction. Players who have just
lived through the night of 1507 deserve to know what actually happened
to the historical protagonists.

\subsubsection{Andrea Doria --- the missing character}

% @ohm:Actor { "id": "genova:andrea-doria",
%   "label": "Andrea Doria",
%   "type": "ohm:Person", "timespan": "1466/1560",
%   "ohm:significance": "Admiral, later Prince of Melfi" }

The real Andrea Doria in 1507 is forty-one years old and already an
accomplished naval commander. He fights for the Fieschi, then the Pope,
then France, then Spain. In \textbf{1528} he switches sides definitively:
he brings Genoa over to the Habsburgs under Charles~V, receives the title
of \textit{Prince of Melfi}, and in practice governs the Republic until
his death in \textbf{1560}. He is the most important figure in
sixteenth-century Genoese history --- and is deliberately absent from
your session because his destiny was already written.

\columnbreak

\subsubsection{The Fieschi --- the failed rebellion}

% @ohm:Event { "id": "genova:congiura-fieschi", "date": "1547",
%   "label": "Fieschi conspiracy against the Doria",
%   "type": "ohm:MilitaryEvent" }

In \textbf{1547} Gian Luigi Fieschi --- grandson of the Niccolò in your
game --- attempts a coup against Andrea Doria. The conspiracy fails by
accident: Fieschi drowned when he fell into the sea from a gangplank on
the night of the coup. The episode inspires Schiller's \textit{Fiesco}
(1783). The Fieschi lost nearly everything after 1547.

\subsubsection{The Spinola --- the continuity of power}

The Spinola survive everything. By the seventeenth and eighteenth centuries
they are among the most powerful families in Europe, with branches in Spain,
Flanders, and the Empire. Their model --- financial capital, diplomatic
network, adaptation to whoever is winning --- proves the most durable.

\subsubsection{The Grimaldi --- Monaco today}

% @ohm:Actor { "id": "genova:grimaldi-monaco",
%   "label": "House Grimaldi of Monaco",
%   "type": "ohm:Dynasty", "timespan": "1297/" }

The House of Grimaldi governs the Principality of Monaco to this day.
This is the direct consequence of the double-game strategy that your
Grimaldi table simulated tonight. Their approach --- maintaining autonomy
by conceding the minimum necessary to the dominant power --- worked for
five hundred years.

\end{multicols}

% =============================================================================
\section{Genoa as Geographic Space: The Sestieri in 1507}
\label{sec:geography}
% =============================================================================

% @ohm:Place { "id": "genova:CittaDiGenova",
%   "label": "Genoa", "type": "ohm:City",
%   "geo:lat": 44.4056, "geo:long": 8.9463,
%   "ohm:accuracy": "high",
%   "ohm:timespan": "ancient/" }

Genoa in 1507 is organised into \textbf{six \textit{sestieri}} (districts),
each with its own social identity, patron churches, and dominant families.
Topography shapes the tactics of the revolt: who controls which
\textit{sestiere} determines who controls the city.

\begin{center}
\rowcolors{2}{LightBG}{white}
\begin{tabularx}{\textwidth}{L{2.5cm} L{3cm} L{2.5cm} X}
  \toprule
  \textbf{Sestiere} & \textbf{Location} & \textbf{Dom.\ Family} & \textbf{Character} \\
  \midrule

  % @ohm:Place { "id": "genova:sestiere-pre",
  %   "label": "Sestiere di Prè", "type": "ohm:District",
  %   "geo:lat": 44.4080, "geo:long": 8.9170,
  %   "ohm:accuracy": "approximate",
  %   "ohm:partOf": "genova:CittaDiGenova" }
  \textbf{Prè} & West, old harbour & Doria & Sailors, arsenals, fishermen. Heart of the popular resistance. \\

  % @ohm:Place { "id": "genova:sestiere-portoria",
  %   "label": "Sestiere di Portoria", "type": "ohm:District",
  %   "geo:lat": 44.4065, "geo:long": 8.9420,
  %   "ohm:accuracy": "approximate" }
  \textbf{Portoria} & Centre-east & Fieschi & Craftsmen, workshops, the Fieschi's popular base. \\

  % @ohm:Place { "id": "genova:sestiere-san-giovanni",
  %   "label": "Sestiere di San Giovanni", "type": "ohm:District",
  %   "geo:lat": 44.4095, "geo:long": 8.9280,
  %   "ohm:accuracy": "approximate" }
  \textbf{San Giovanni} & Centre, harbour & Spinola & Merchants, bankers, the financial heart of the city. \\

  % @ohm:Place { "id": "genova:sestiere-san-lorenzo",
  %   "label": "Sestiere di San Lorenzo", "type": "ohm:District",
  %   "geo:lat": 44.4071, "geo:long": 8.9320,
  %   "ohm:accuracy": "approximate" }
  \textbf{San Lorenzo} & Historic centre & Fieschi & The cathedral, the ducal palace, civic and religious power. \\

  % @ohm:Place { "id": "genova:sestiere-castello",
  %   "label": "Sestiere di Castello", "type": "ohm:District",
  %   "geo:lat": 44.4060, "geo:long": 8.9380,
  %   "ohm:accuracy": "approximate" }
  \textbf{Castello} & East, hillside & Doria/Spinola & Mixed nobility, palaces, defensible elevated position. \\

  % @ohm:Place { "id": "genova:sestiere-soziglia",
  %   "label": "Sestiere di Soziglia", "type": "ohm:District",
  %   "geo:lat": 44.4082, "geo:long": 8.9350,
  %   "ohm:accuracy": "approximate" }
  \textbf{Soziglia} & Centre-north & Grimaldi & International trade, foreign \textit{fondachi}, transit zone. \\
  \bottomrule
\end{tabularx}
\end{center}

\begin{multicols}{2}

\subsection{Key Locations of the Revolt}

% @ohm:Place { "id": "genova:SanLorenzo",
%   "label": "Cathedral of San Lorenzo",
%   "type": "ohm:ReligiousBuilding",
%   "geo:lat": 44.4071, "geo:long": 8.9316,
%   "ohm:accuracy": "high",
%   "ohm:significance": "Spiritual and political centre of the revolt" }

\textbf{Cathedral of San Lorenzo} --- The bells of San Lorenzo are the
revolt's signal. The cathedral is also where informal civic assemblies
are held. Whoever controls San Lorenzo controls the city's \textit{symbol}.

% @ohm:Place { "id": "genova:PalazzoDucale",
%   "label": "Ducal Palace (Palazzo Ducale)",
%   "type": "ohm:GovernmentBuilding",
%   "geo:lat": 44.4068, "geo:long": 8.9337,
%   "ohm:accuracy": "high" }

\textbf{Ducal Palace} --- Seat of formal government. In 1507 it is held
by the French administration. Whoever controls it at the revolt's end
holds executive power.

\columnbreak

% @ohm:Place { "id": "genova:Porto",
%   "label": "Harbour of Genoa",
%   "type": "ohm:Harbour",
%   "geo:lat": 44.4063, "geo:long": 8.9237,
%   "ohm:accuracy": "high" }

\textbf{The Harbour} --- The economic heart. The Doria's galleys are
anchored here. Whoever controls the harbour controls supplies and
escape routes --- for the French and for defeated rebels alike.

% @ohm:Place { "id": "genova:Molo",
%   "label": "Il Molo (main pier)",
%   "type": "ohm:Infrastructure",
%   "geo:lat": 44.4046, "geo:long": 8.9268,
%   "ohm:accuracy": "approximate" }

\textbf{Il Molo} --- The main pier, control point for all goods entering
and leaving. Taxed by the French, guarded by a mixed garrison.

\end{multicols}

% =============================================================================
\section{Genoa as a Public History Event}
\label{sec:public-history}
% =============================================================================

\begin{rulebox}[GoldRPG]{For Organisers and Facilitators}
This game is designed to function as a \textbf{public history event} ---
a format in which ludic participation becomes a vehicle for historical
understanding. Practical guidance:

\medskip
\textbf{Before the session (30 minutes):}
Briefly present the historical context using the chronology in
\S\ref{sec:revolt-chronology}. Players do not need to know every detail
--- the goal is a \textit{sense of authenticity} that makes their choices
feel meaningful.

\medskip
\textbf{During the session:}
GMs may introduce historical details as in-game information: an NPC
mentioning the commercial losses after the fall of Caffa, a letter citing
Louis~XII's taxes. Fiction and history reinforce each other.

\medskip
\textbf{After the session (45 minutes):}
Structured debrief in two parts:
\begin{enumerate}
  \item \textit{What happened in the game?} --- each table narrates
        its own epilogue.
  \item \textit{What actually happened?} --- comparison with the real
        chronology (\S\ref{sec:aftermath}). Players discover what happened
        to ``their'' historical figures in the following decades.
\end{enumerate}

\end{rulebox}
